\documentclass[11pt,a4paper]{article}

\usepackage{fontspec}
\usepackage{polyglossia}
\setmainlanguage{english}
\setotherlanguage{russian}
\setmainfont{CMU Serif}
\setsansfont{CMU Sans Serif}
\setmonofont{CMU Typewriter Text}

\usepackage{amsmath,amssymb,amsthm,mathtools}
\usepackage{booktabs}
\usepackage{enumitem}
\usepackage{geometry}
\usepackage{hyperref}
\usepackage{microtype}

\geometry{margin=2.5cm}
\setlength{\parindent}{1.2em}
\setlength{\parskip}{0.35em}
\sloppy

\hypersetup{
  colorlinks=true,
  linkcolor=black,
  citecolor=black,
  urlcolor=blue,
  pdftitle={Unified Whole Theory: A Mereological and Relational Foundation for Space-Time},
  pdfauthor={Andrey Tikhonov (Тихонов Андрей)}
}

\theoremstyle{definition}
\newtheorem{definition}{Definition}[section]
\newtheorem{axiom}{Axiom}[section]
\newtheorem{postulate}{Postulate}[section]
\newtheorem{definitionru}{Определение}[section]
\newtheorem{axiomru}{Аксиома}[section]
\newtheorem{postulateru}{Постулат}[section]

\theoremstyle{plain}
\newtheorem{theorem}{Theorem}[section]
\newtheorem{proposition}{Proposition}[section]
\newtheorem{corollary}{Corollary}[section]
\newtheorem{theoremru}{Теорема}[section]
\newtheorem{propositionru}{Предложение}[section]
\newtheorem{corollaryru}{Следствие}[section]

\theoremstyle{remark}
\newtheorem{remark}{Remark}[section]
\newtheorem{remarkru}{Замечание}[section]

\newcommand{\Whole}{U}
\newcommand{\Parts}{\mathcal A}
\newcommand{\Rel}{R}
\newcommand{\Values}{V}
\newcommand{\Vstruct}{\mathbf V}
\newcommand{\State}{S}
\newcommand{\States}{\mathcal S}
\newcommand{\PartOf}{\sqsubseteq}
\newcommand{\PropPart}{\sqsubset}
\newcommand{\comp}{\circ}
\newcommand{\Vord}{\preceq}
\newcommand{\norm}{\rho}
\newcommand{\Disc}{\mathsf{Disc}}
\newcommand{\Space}{\mathsf{Space}}
\newcommand{\Time}{\mathsf{Time}}
\newcommand{\Motion}{\mathsf{Motion}}
\newcommand{\Observer}{\mathsf{Obs}}
\newcommand{\Struct}{\operatorname{Struct}}
\newcommand{\Inv}{\operatorname{Inv}}
\newcommand{\Stab}{\sigma}
\newcommand{\Cost}{\mathcal C}

\title{\textbf{Unified Whole Theory}\\
\large A Mereological and Relational Foundation for Space-Time}
\author{Andrey Tikhonov (Тихонов Андрей)\\\texttt{uwt@xteam.pro}}
\date{July 10, 2026}

\begin{document}
\maketitle

\begin{abstract}
Unified Whole Theory (UWT) is proposed as a formal relational framework in which
space, time, motion, distance, mass, and energy are not primitive notions but
derived structures. The construction starts from a whole \(\Whole\), a
mereological lattice of parts \(\Parts\), distinguishability, and relations
valued in a normalized ordered group \(\Vstruct\). Relational distance is
obtained from normalized relation values, the triangle inequality is derived
from relation composition rather than postulated, time is introduced as a
structure of changes between states, and elementary mechanical quantities are
defined as invariants of relational change. The paper gives an explicit finite
model establishing relative consistency of the basic axioms and lists research
directions where this relations-first framework may differ from conventional
space-time-first approaches. The Russian version follows the English version in
full.
\end{abstract}

\tableofcontents

\section{Introduction}

Most physical theories begin with a stage: space, time, coordinates, distances,
and often a differentiable or metric structure. Objects are then placed on that
stage and their evolution is described. UWT reverses this order. Its guiding
claim is
\[
\boxed{\text{relations are not in space and time; space and time arise from relations.}}
\]
The aim is not to deny the empirical usefulness of space-time models, but to
reconstruct the conceptual order in which such models become derived structures.

The English and Russian versions in this submission are intentionally contained
in one document. This follows the arXiv rule that multilingual submissions must
provide a complete English version first, followed by the non-English version.

\section{Basic ontology: whole, parts, and distinguishability}

\begin{postulate}[The whole is not externally placed]
The whole \(\Whole\) is not an object inside a prior space or a prior time. If
space and time occur, they occur as internal structures of \(\Whole\).
\end{postulate}

\begin{definition}[Parts]
\(\Parts\) denotes the domain of distinguishable parts of \(\Whole\). A relation
\(A\PartOf B\) reads ``\(A\) is part of \(B\)''.
\end{definition}

\begin{axiom}[Mereological bounded lattice]
\((\Parts,\PartOf)\) is a bounded lattice with bottom \(0\) and top \(\Whole\).
Thus every pair of parts has a meet and a join, and all parts are contained in
\(\Whole\).
\end{axiom}

\begin{definition}[Distinguishability]
\[
\Disc(A_i,A_j) \iff A_i \neq A_j .
\]
Distinguishability is prior to metric distance. If parts are not
distinguishable, no relation that could later be normalized as distance can be
specified.
\end{definition}

\section{Relation values}

\begin{definition}[Normalized ordered relation group]
A relation-value structure is a tuple
\[
\Vstruct=(\Values,\comp,e,\preceq,\norm),
\]
where \((\Values,\comp,e)\) is a group, \(\preceq\) is an order compatible with
composition, and \(\norm:\Values\to\mathbb R_{\ge 0}\) is a non-negative
normalization satisfying \(\norm(e)=0\) and
\[
\norm(v\comp w)\leq \norm(v)+\norm(w).
\]
\end{definition}

\begin{definition}[Relation]
A relation is a map
\[
\Rel:\Parts\times\Parts\to\Values .
\]
The value \(\Rel(A_i,A_j)\) is not assumed to be a number. Numerical quantities
appear only after applying the normalization \(\norm\).
\end{definition}

\begin{remark}
This separation prevents hidden assumptions about coordinates or metric
backgrounds. A relation can carry order, composition, orientation, or other
structural content before any numerical projection is chosen.
\end{remark}

\section{Space as a derived metric structure}

\begin{definition}[Relational distance]
\[
d(A_i,A_j)=\norm(\Rel(A_i,A_j)).
\]
\end{definition}

\begin{theorem}[Triangle inequality]
Assume relation composition satisfies
\[
\Rel(A_i,A_k)=\Rel(A_i,A_j)\comp\Rel(A_j,A_k)
\]
for composable triples and the normalization is subadditive. Then
\[
d(A_i,A_k)\leq d(A_i,A_j)+d(A_j,A_k).
\]
\end{theorem}

\begin{proof}
By definition and subadditivity,
\[
d(A_i,A_k)=\norm(\Rel(A_i,A_k))
=\norm(\Rel(A_i,A_j)\comp\Rel(A_j,A_k))
\leq \norm(\Rel(A_i,A_j))+\norm(\Rel(A_j,A_k)).
\]
The right hand side is \(d(A_i,A_j)+d(A_j,A_k)\).
\end{proof}

Thus a metric property normally postulated in geometric theories appears here as
a theorem about relation composition and normalization. ``Space'' is the
invariant structure generated by these distances, not a primitive container.

\section{States, change, and relational time}

\begin{definition}[State]
A state \(\State\in\States\) is a complete relational configuration over the
parts of \(\Whole\). Formally it may be represented as a map
\[
\State:\Parts\times\Parts\to\Values .
\]
\end{definition}

\begin{definition}[Relational change]
For two states \(S_1,S_2\), their change is the structured difference
\[
\Delta(S_1,S_2)=\{\Rel_{S_1}(A_i,A_j)^{-1}\comp \Rel_{S_2}(A_i,A_j)\}_{i,j}.
\]
\end{definition}

\begin{definition}[Relational time]
Time is a measure or order of changes between states:
\[
\Time(S_1,S_2)=T(\Delta(S_1,S_2)).
\]
There is no external clock at the foundational level; clocks are stable
subsystems whose internal relational changes can be used as reference processes.
\end{definition}

\section{Observers and relative centers}

An observer is not an entity outside the whole. It is a stable subsystem
\(O\in\Parts\) whose relations to other parts define a relative center and a
reference order. Observed space and observed time are therefore perspective
structures built from relational invariants around \(O\), not absolute
containers independent of all subsystems.

\section{Relational mechanics}

\begin{definition}[Velocity and acceleration]
Given a relational path \(A(t)\), velocity and acceleration are defined by
changes of relation values with respect to relational time:
\[
v=\frac{d}{dT}\Rel(A(t_1),A(t_2)),\qquad
a=\frac{d v}{dT}.
\]
The notation is shorthand for a limit or discrete approximation inside a chosen
state-transition structure.
\end{definition}

\begin{definition}[Stability and mass]
The stability of a part is an invariant \(\Stab(A)\) of its relations under
state changes. Mass is modeled as a monotone function of stability,
\[
m=M(\Stab(A)),
\]
or, equivalently in a local dynamical regime, as curvature of a cost functional.
\end{definition}

\begin{definition}[Cost of change]
Let \(\Cost(v)\) be the cost of relational change. Around a stable state,
\[
\Cost(v)=\Cost(0)+\frac12\Cost''(0)v^2+O(v^3).
\]
Identifying \(m=\Cost''(0)\) gives the kinetic form
\[
E_k=\frac12 m v^2
\]
as a derived approximation rather than a primitive postulate.
\end{definition}

\section{Relative consistency}

\begin{theorem}[Finite model]
The basic axioms are relatively consistent.
\end{theorem}

\begin{proof}[Model sketch]
Let \(\Parts\) be the finite powerset lattice of a non-empty finite set with
inclusion as \(\PartOf\). Let \(\Values=(\mathbb R,+,0,\leq,|\cdot|)\). Define
\(\Rel(A,B)=\mu(A\triangle B)\), where \(\mu\) is cardinality or a normalized
finite measure and \(A\triangle B\) is symmetric difference. Composition is
addition along composable transitions, and the norm is absolute value. The
lattice axioms, identity, ordering, and subadditivity are satisfied in this
model. Therefore the listed axioms do not entail contradiction relative to
ordinary finite set theory and real arithmetic.
\end{proof}

\section{Empirical and mathematical program}

UWT becomes physically useful only when it yields models that reproduce known
limits and produce distinguishing predictions. The immediate program is:
\begin{enumerate}[label=(\roman*)]
  \item classify admissible relation-value structures \(\Vstruct\);
  \item derive metric, curvature, and dimensionality from relation composition;
  \item identify observer subsystems and clock subsystems;
  \item recover classical mechanics and relativistic limits where appropriate;
  \item compare relational wavefunction models with quantum-mechanical
        dispersion, stability, and measurement behavior.
\end{enumerate}

\section{Conclusion}

The central result of this formulation is an ordering of concepts:
\[
\Disc \Rightarrow \Rel \Rightarrow \State \Rightarrow \Space,\Time
\Rightarrow \Motion \Rightarrow \text{physics}.
\]
Coordinates, distance, velocity, mass, and energy are not removed; they are
reintroduced as derived relational structures. The approach is therefore a
foundation program rather than a replacement of existing empirical theories at
the first step.

\newpage
\begin{russian}

\section*{Русская версия}
\addcontentsline{toc}{section}{Русская версия}

\begin{center}
{\Large\textbf{Теория Единого Целого}}\\
\large Мереологическое и реляционное основание пространства-времени
\end{center}

\begin{abstract}
Теория Единого Целого (ТЕЦ) предлагается как формальный реляционный каркас, в
котором пространство, время, движение, расстояние, масса и энергия не являются
первичными понятиями, а выводятся как производные структуры. Построение
начинается с целого \(\Whole\), мереологической решётки частей \(\Parts\),
различимости и отношений со значениями в нормированной упорядоченной группе
\(\Vstruct\). Реляционное расстояние получается из нормированных значений
отношений, неравенство треугольника выводится из композиции отношений, время
задаётся как структура изменений между состояниями, а механические величины
определяются как инварианты реляционного изменения. Также приводится конечная
модель, устанавливающая относительную непротиворечивость базовых аксиом.
\end{abstract}

\section{Введение}

Большинство физических теорий начинают с готовой сцены: пространства, времени,
координат, расстояний и часто дифференцируемой или метрической структуры. Затем
на эту сцену помещаются объекты и описывается их эволюция. ТЕЦ меняет порядок:
\[
\boxed{\text{не отношения существуют в пространстве и времени,}}
\]
\[
\boxed{\text{а пространство и время возникают из отношений.}}
\]
Задача теории не в отрицании практической ценности пространственно-временных
моделей, а в восстановлении порядка вывода, при котором такие модели становятся
производными структурами.

\section{Базовая онтология: целое, части и различимость}

\begin{postulateru}[Целое не помещено вовне]
Целое \(\Whole\) не является объектом внутри заранее данного пространства или
заранее данного времени. Если пространство и время возникают, они возникают как
внутренние структуры \(\Whole\).
\end{postulateru}

\begin{definitionru}[Части]
\(\Parts\) обозначает область различимых частей \(\Whole\). Запись
\(A\PartOf B\) читается как «\(A\) является частью \(B\)».
\end{definitionru}

\begin{axiomru}[Ограниченная мереологическая решётка]
\((\Parts,\PartOf)\) является ограниченной решёткой с нижним элементом \(0\) и
верхним элементом \(\Whole\). Поэтому любая пара частей имеет пересечение и
объединение, а все части содержатся в \(\Whole\).
\end{axiomru}

\begin{definitionru}[Различимость]
\[
\Disc(A_i,A_j)\iff A_i\neq A_j.
\]
Различимость предшествует метрическому расстоянию. Если части не различимы,
нельзя задать отношение, которое позднее могло бы быть нормировано как
расстояние.
\end{definitionru}

\section{Значения отношений}

\begin{definitionru}[Нормированная упорядоченная группа значений]
Структура значений отношений есть кортеж
\[
\Vstruct=(\Values,\comp,e,\preceq,\norm),
\]
где \((\Values,\comp,e)\) — группа, \(\preceq\) — порядок, совместимый с
композицией, а \(\norm:\Values\to\mathbb R_{\ge0}\) — неотрицательная
нормировка, удовлетворяющая \(\norm(e)=0\) и
\[
\norm(v\comp w)\leq \norm(v)+\norm(w).
\]
\end{definitionru}

\begin{definitionru}[Отношение]
Отношение есть отображение
\[
\Rel:\Parts\times\Parts\to\Values.
\]
Значение \(\Rel(A_i,A_j)\) не предполагается числом. Числовые величины
появляются только после применения нормировки \(\norm\).
\end{definitionru}

\section{Пространство как производная метрическая структура}

\begin{definitionru}[Реляционное расстояние]
\[
d(A_i,A_j)=\norm(\Rel(A_i,A_j)).
\]
\end{definitionru}

\begin{theoremru}[Неравенство треугольника]
Пусть композиция отношений удовлетворяет
\[
\Rel(A_i,A_k)=\Rel(A_i,A_j)\comp\Rel(A_j,A_k)
\]
для компонуемых троек, а нормировка субаддитивна. Тогда
\[
d(A_i,A_k)\leq d(A_i,A_j)+d(A_j,A_k).
\]
\end{theoremru}

\begin{proof}
По определению и субаддитивности,
\[
d(A_i,A_k)=\norm(\Rel(A_i,A_k))
=\norm(\Rel(A_i,A_j)\comp\Rel(A_j,A_k))
\leq \norm(\Rel(A_i,A_j))+\norm(\Rel(A_j,A_k)).
\]
Правая часть равна \(d(A_i,A_j)+d(A_j,A_k)\).
\end{proof}

Тем самым метрическое свойство, обычно постулируемое в геометрических теориях,
становится теоремой о композиции и нормировке отношений.

\section{Состояния, изменение и реляционное время}

\begin{definitionru}[Состояние]
Состояние \(\State\in\States\) — это полная реляционная конфигурация частей
\(\Whole\). Формально его можно представить как отображение
\[
\State:\Parts\times\Parts\to\Values.
\]
\end{definitionru}

\begin{definitionru}[Реляционное изменение]
Для двух состояний \(S_1,S_2\) их изменение задаётся структурной разностью
\[
\Delta(S_1,S_2)=\{\Rel_{S_1}(A_i,A_j)^{-1}\comp \Rel_{S_2}(A_i,A_j)\}_{i,j}.
\]
\end{definitionru}

\begin{definitionru}[Реляционное время]
Время есть мера или порядок изменений между состояниями:
\[
\Time(S_1,S_2)=T(\Delta(S_1,S_2)).
\]
В основании нет внешних часов; часы — это устойчивые подсистемы, внутреннее
реляционное изменение которых используется как опорный процесс.
\end{definitionru}

\section{Наблюдатели и относительные центры}

Наблюдатель не находится вне целого. Это устойчивая подсистема \(O\in\Parts\),
отношения которой с другими частями задают относительный центр и порядок
отсчёта. Наблюдаемое пространство и наблюдаемое время являются перспективными
структурами, построенными из реляционных инвариантов вокруг \(O\).

\section{Реляционная механика}

\begin{definitionru}[Скорость и ускорение]
Для реляционного пути \(A(t)\) скорость и ускорение определяются изменениями
значений отношений относительно реляционного времени:
\[
v=\frac{d}{dT}\Rel(A(t_1),A(t_2)),\qquad
a=\frac{dv}{dT}.
\]
\end{definitionru}

\begin{definitionru}[Устойчивость и масса]
Устойчивость части — это инвариант \(\Stab(A)\) её отношений при изменениях
состояния. Масса моделируется как монотонная функция устойчивости,
\[
m=M(\Stab(A)),
\]
или в локальном динамическом режиме как кривизна функционала стоимости.
\end{definitionru}

\begin{definitionru}[Стоимость изменения]
Пусть \(\Cost(v)\) — стоимость реляционного изменения. Вблизи устойчивого
состояния
\[
\Cost(v)=\Cost(0)+\frac12\Cost''(0)v^2+O(v^3).
\]
Отождествление \(m=\Cost''(0)\) даёт кинетическую форму
\[
E_k=\frac12mv^2
\]
как производное приближение, а не первичный постулат.
\end{definitionru}

\section{Относительная непротиворечивость}

\begin{theoremru}[Конечная модель]
Базовые аксиомы относительно непротиворечивы.
\end{theoremru}

\begin{proof}
Пусть \(\Parts\) — конечная решётка подмножеств непустого конечного множества с
включением как \(\PartOf\). Пусть
\(\Values=(\mathbb R,+,0,\leq,|\cdot|)\). Определим
\(\Rel(A,B)=\mu(A\triangle B)\), где \(\mu\) — мощность или нормированная
конечная мера, а \(A\triangle B\) — симметрическая разность. Композиция
задаётся сложением вдоль компонуемых переходов, нормировка — модулем. В этой
модели выполняются аксиомы решётки, тождества, порядка и субаддитивности.
\end{proof}

\section{Программа применения}

Физическая содержательность ТЕЦ требует воспроизведения известных пределов и
получения различающих следствий. Ближайшая программа:
\begin{enumerate}[label=(\roman*)]
  \item классифицировать допустимые структуры значений \(\Vstruct\);
  \item вывести метрику, кривизну и размерность из композиции отношений;
  \item определить подсистемы-наблюдатели и подсистемы-часы;
  \item восстановить классическую механику и релятивистские пределы;
  \item сопоставить реляционные модели волновой функции с дисперсией,
        устойчивостью и поведением измерения в квантовой механике.
\end{enumerate}

\section{Заключение}

Ключевой порядок понятий:
\[
\Disc \Rightarrow \Rel \Rightarrow \State \Rightarrow \Space,\Time
\Rightarrow \Motion \Rightarrow \text{физика}.
\]
Координаты, расстояние, скорость, масса и энергия не исчезают; они возвращаются
как производные реляционные структуры.

\end{russian}

\begin{thebibliography}{9}

\bibitem{rovelli1996}
C. Rovelli,
``Relational quantum mechanics,''
\emph{International Journal of Theoretical Physics} 35, 1637--1678 (1996).

\bibitem{barbour1999}
J. Barbour,
\emph{The End of Time: The Next Revolution in Physics},
Oxford University Press, 1999.

\bibitem{baez2006}
J. C. Baez and J. Dolan,
``From finite sets to Feynman diagrams,''
in \emph{Mathematics Unlimited---2001 and Beyond}, Springer, 2001.

\bibitem{szabo2015}
M. E. Szabo,
\emph{The Collected Papers of Gerhard Gentzen},
North-Holland, 1969.

\end{thebibliography}

\end{document}
